\chapter{Lieux}\label{ch:in:locations}

Un package de 300\,000~\$, après impôt sur le revenu et un an de loyer d'un deux-pièces selon la statistique officielle de chaque ville, laisse
139\,220~\$ à Londres et 214\,938~\$ à Hong Kong. L'ordre n'est pas figé~: New York laisse plus que Londres à ce package et moins
au-dessus d'environ 851\,900~\$, parce que les deux barèmes d'impôt se croisent et que les loyers sont presque égaux. Ce chapitre
décrit les villes de l'industrie, calcule ce que vaut un package après impôt et logement dans cinq d'entre elles, présente les
visas qu'un déménagement exige et discute ce que coûte un changement de ville.

\section{Les villes de l'industrie et ce qu'on y trouve}

Les employeurs de l'industrie se concentrent là où se trouvent ses marchés et ses clients (chapitre~1)~: New York et Chicago pour les actions américaines,
les options et les contrats à terme~; Londres pour les devises, les taux et les marchés européens~; Amsterdam pour les options européennes et le market making
sur ETF (chapitres~2 et~3)~; Zurich pour les banques~; Hong Kong et Singapour pour l'Asie~; et Dubaï, dont l'absence d'impôt sur le revenu
que mesure le chapitre~15. La colocation (livre~14, chapitre~9) rattache les systèmes de trading aux
centres de données des bourses, mais pas les personnes~: les chercheurs et les ingénieurs d'une société peuvent travailler loin de ses serveurs. Le
chapitre compare cinq villes dont les statistiques officielles de loyers sont publiées~: Londres, New York, Chicago, Zurich et Hong Kong.

\begin{definition}[Indice de niveau des prix, rémunération disponible après logement]\label{def:in:locations:pli}
Un \emph{indice de niveau des prix}\index{indice de niveau des prix} est le rapport entre la parité de pouvoir d'achat d'un pays et son taux de change de marché,
exprimé par rapport à une zone de référence~: au-dessus de 100, le même panier coûte plus cher là-bas que dans la référence.
\emph{rémunération disponible après logement}\index{rémunération disponible après logement} est la rémunération d'un package après impôt sur le revenu et cotisations
sociales salariées, diminuée d'un an de loyer d'un logement standard de la ville.
\end{definition}

\section{Rémunération, impôt et logement}

\begin{dated}{2026-08}{Loyers dans les cinq villes du chapitre}\label{dat:in:locations:rents}
Londres, août 2026~: loyer privé moyen d'un deux-pièces dans la région de Londres, 2\,232~\pounds par mois (ONS).
New York et Chicago, exercice fiscal 2026~: loyers de marché de référence d'un deux-pièces, 2\,910~\$ et 1\,781~\$ par mois~; les
loyers de marché de référence «~sont des estimations des loyers bruts du 40\up{e}~centile pour des logements de qualité standard~» (HUD). Zurich, mars 2024~: loyer net médian
d'un appartement de trois pièces de la ville, 1\,578~francs par mois (City of Zurich). Hong Kong, août 2026 (provisoire)~:
loyer moyen des appartements de classe~B (40 à 69.9 mètres carrés) sur Hong Kong Island, 446~HK\$ par mètre carré et par mois (Rating and
Valuation Department). \omterm{def:in:locations:pli}{Indices de niveau des prix} de la consommation individuelle effective, 2024 (UE = 100)~: Suisse 184.3, États-Unis
149.2, Royaume-Uni 129.1, Pays-Bas 121.0 (Eurostat).
\end{dated}

Les cinq statistiques mesurent des choses différentes~: une moyenne, un 40\up{e}~centile, une médiane~; des loyers charges comprises
(loyers bruts de HUD) et hors charges (loyer net de Zurich)~; une région (Londres), une aire métropolitaine (New York, Chicago) et une ville
(Zurich). Pour Hong Kong, le chapitre multiplie le loyer au mètre carré par une surface illustrative de 55~mètres carrés. Le
comparaison relève donc d'ordres de grandeur et de la manière dont le \omterm{def:in:bank-markets-divisions:league}{classement} se déplace, non d'écarts exacts.

\begin{method}[Rémunération disponible après logement]\label{met:in:locations:disposable}
Pour un package $G$ en dollars américains et une ville $c$ de fonction fiscale $T_c$ (chapitre~15, \texttt{firm.aftertax}, aux
taux de change moyens 2025 de la BCE) et loyer mensuel $r_c$ en monnaie locale avec $k_c$ unités par dollar,
\[
D_c(G)=G-T_c(G)-12\,r_c/k_c .
\]
Deux villes se comparent par le package $G_b$ tel que $D_b(G_b)=D_a(G_a)$, et par les packages pour lesquels $D_a=D_b$.
\end{method}

\begin{omfigure}
\begin{tikzpicture}
\begin{axis}[xbar stacked,width=10.5cm,height=5cm,xmin=0,xmax=300,ymin=0.4,ymax=5.6,bar width=11pt,ytick=data,
  yticklabels={{Hong Kong},{Zurich},{Chicago},{New York},{Londres}},yticklabel style={font=\footnotesize},
  xticklabel style={font=\footnotesize},xlabel={\small milliers de \$ pour un package de 300\,000~\$},table/col sep=comma,
  legend style={legend cell align=left,font=\scriptsize,at={(0.5,-0.32)},anchor=north,legend columns=3,/tikz/every even column/.append style={column sep=8pt}}]
\addplot[fill=omDef,draw=none] table[x=disposable,y=pos]{figdata/industry/27-locations/split.csv};
\addplot[fill=omThm!60,draw=none] table[x=rent,y=pos]{figdata/industry/27-locations/split.csv};
\addplot[fill=black!25,draw=none] table[x=tax,y=pos]{figdata/industry/27-locations/split.csv};
\legend{rémunération disponible après logement,loyer,impôt et cotisations sociales}
\end{axis}
\end{tikzpicture}

\omcaption{Un package de 300\,000~\$ réparti entre la \omterm{def:in:locations:pli}{rémunération disponible après logement}, un an de loyer d'un deux-pièces selon
la statistique officielle de chaque ville, et impôt avec cotisations sociales salariées (salarié seul, sans autre revenu, déductions standard).
Données~: \texttt{data/industry/rents.csv}, données fiscales du chapitre~15 et deux constantes, via
\texttt{in\_locations.table}.}\label{fig:in:locations:split}
\end{omfigure}

Pour un package de 300\,000~\$ (\cref{fig:in:locations:split}), l'ordre est Hong Kong 214\,938~\$, Zurich 184\,855~\$, Chicago 178\,955~\$,
New York 149\,062~\$ et Londres 139\,220~\$. L'impôt sépare davantage les villes que le loyer~: l'impôt et les cotisations de Londres
prélèvent 125\,455~\$ et ceux de Hong Kong 47\,309~\$, tandis que les loyers vont de 21\,372~\$ (Chicago) à 37\,753~\$ (Hong Kong). Le
L'ordre évolue avec le package (\cref{fig:in:locations:vslondon})~: New York devance Londres d'environ 134\,100~\$
à environ 851\,900~\$ et reste derrière Londres en dehors de cette plage, parce que le barème new-yorkais taxe plus lourdement les gros revenus que
le barème londonien (chapitre~15), alors que leurs loyers se situent à moins de 400~\$ par an l'un de l'autre.

\begin{omfigure}
\begin{tikzpicture}
\begin{axis}[width=10.5cm,height=5.8cm,xmin=100,xmax=1000,ymin=-15,ymax=80,
  xlabel={\small package, milliers de \$},ylabel={\small écart avec Londres, milliers de \$},
  tick label style={font=\footnotesize},grid=major,grid style={gray!20},table/col sep=comma,
  legend style={font=\scriptsize,at={(0.97,0.22)},anchor=south east,legend cell align=left}]
\addplot[omThm,thick] table[x=gross,y=Hong_Kong]{figdata/industry/27-locations/vs_london.csv};
\addplot[black!70,thick,dashed] table[x=gross,y=Zurich]{figdata/industry/27-locations/vs_london.csv};
\addplot[omDef!60,thick] table[x=gross,y=Chicago]{figdata/industry/27-locations/vs_london.csv};
\addplot[omDef,thick] table[x=gross,y=New_York]{figdata/industry/27-locations/vs_london.csv};
\draw[black!50] (axis cs:100,0) -- (axis cs:1000,0);
\legend{Hong Kong,Zurich,Chicago,New York}
\end{axis}
\end{tikzpicture}

\omcaption{\omterm{def:in:locations:pli}{Rémunération disponible après logement} dans quatre villes, moins celle de Londres, pour des packages de 100\,000~\$ à 1~million de dollars. La
courbe de New York coupe zéro deux fois~; la courbe de Hong Kong sort du graphique juste au-dessus de 300\,000~\$ et atteint 296\,131~\$ pour 1
million de dollars. Données~: \texttt{in\_locations.curves}.}\label{fig:in:locations:vslondon}
\end{omfigure}

Les niveaux de prix ajoutent une troisième couche que le chapitre ne modélise pas~: la même rémunération disponible permet d'acheter moins en Suisse,
dont le niveau de prix de la consommation des ménages valait 184.3 contre 100 pour l'Union européenne en 2024, que
celui du Royaume-Uni (129.1) ou des Pays-Bas (121.0) (\cref{fig:in:locations:pli}). Le chapitre déduit le loyer parce que les
statistiques officielles le mesurent ville par ville~; le reste du coût de la vie se traite mieux avec un tel indice que par
des estimations propres à chaque ville.

\begin{omfigure}
\begin{tikzpicture}
\begin{axis}[xbar,width=10.5cm,height=4cm,xmin=0,xmax=210,ymin=0.4,ymax=4.6,bar width=9pt,ytick=data,
  yticklabels={{Pays-Bas},{Royaume-Uni},{États-Unis},{Suisse}},yticklabel style={font=\footnotesize},
  xticklabel style={font=\footnotesize},xlabel={\small indice de niveau des prix, consommation individuelle effective, UE = 100, 2024},
  grid=major,grid style={gray!20},table/col sep=comma,nodes near coords,
  nodes near coords style={font=\scriptsize,anchor=west,xshift=1pt},point meta=explicit symbolic]
\addplot[fill=omDef,draw=none] table[x=pli,y=pos,meta=pli]{figdata/industry/27-locations/pli.csv};
\end{axis}
\end{tikzpicture}

\omcaption{\omterm{def:in:locations:pli}{Indices de niveau des prix} de la consommation des ménages dans quatre pays, 2024, issus des parités de pouvoir
d'achat d'Eurostat. Données~: Eurostat \texttt{prc\_ppp\_ind}, via \texttt{in\_locations.PLI\_2024}.}\label{fig:in:locations:pli}
\end{omfigure}

\section{Visas et permis de travail}

\begin{definition}[Visa pour profession spécialisée, visa de travailleur qualifié]\label{def:in:locations:visas}
Un \emph{visa pour profession spécialisée}\index{visa pour profession spécialisée} est un statut de travail temporaire aux États-Unis (H-1B) pour des emplois
exigeant au moins un diplôme de premier cycle dans un domaine précis, parrainé par un employeur qui dépose une demande de
condition de travail (chapitre~14). Un \emph{visa de travailleur qualifié}\index{visa de travailleur qualifié} est la voie de travail parrainée du
Royaume-Uni, pour un emploi requérant un niveau de compétence donné et rémunéré à au moins un seuil de salaire.
\end{definition}

\begin{dated}{2026-09}{Visas pour les villes de l'industrie}\label{dat:in:locations:visas}
États-Unis (USCIS, septembre 2026)~: plafond ordinaire H-1B de 65\,000 par an, plus 20\,000 pour les titulaires d'un master américain~;
à partir de l'exercice fiscal 2027, une «~procédure de sélection pondérée qui favorise l'attribution des visas H-1B aux candidats
plus qualifiés et mieux rémunérés~»~; une proclamation du 19 septembre 2025 exige un versement supplémentaire de 100\,000~\$ pour certaines
demandes déposées à partir du 21 septembre 2025, dont les modalités d'application ont été annulées par un tribunal de district le 8 juin 2026, une ordonnance
administrativement suspendue pendant que le gouvernement fait appel. Royaume-Uni~: \omterm{def:in:locations:visas}{visa de travailleur qualifié} «~à au moins 41\,700~\pounds
par an, ou le «~taux courant~» de l'emploi, le plus élevé des deux~». Singapour~: salaire minimal qualifiant pour l'Employment Pass dans les
services financiers 6\,200~S\$ par mois, qui augmente avec l'âge jusqu'à 11\,800~S\$, et 6\,600~S\$ pour les nouvelles demandes à partir du 1\up{er} janvier
2027, plus le test de points COMPASS. Pays-Bas, 2026~: migrant hautement qualifié ou carte bleue européenne, 5\,942~EUR brut par mois
(30 ans et plus), 4\,357~EUR (moins de 30 ans).
\end{dated}

Pour une recrue en quantitative finance, les seuils ne sont guère contraignants (les offres du chapitre~14 se situent bien au-dessus)~; les vraies contraintes se trouvent
ailleurs. La loterie américaine détermine si un premier H-1B est possible une année donnée, et la pondération de 2027 ainsi que
le versement contesté de 100\,000~\$ modifient les chances et le coût pour l'employeur. Les autres voies nationales relèvent davantage
d'un contrôle administratif, ce qui est l'une des raisons pour lesquelles les entreprises ouvrent des bureaux là où l'embauche depuis l'étranger est prévisible.

\section{Déménager~: relocalisation, travail à distance et coût du départ}

Un déménagement comporte trois coûts, outre l'impôt et le loyer de la nouvelle ville. \textbf{La relocalisation}~: billets d'avion, logement temporaire, cautions et
les mois de double loyer, que l'employeur peut prendre en charge pour une recrue. \textbf{Le coût du départ}~: la rémunération différée non acquise et les clauses de
non-concurrence chez l'ancien employeur (chapitres~13 et~28), que le nouvel employeur peut racheter. \textbf{La
\omterm{def:in:pay-regulation-and-tax-by-location:expat}{résidence fiscale}}~: un déménagement en cours d'année peut partager une même année entre deux régimes fiscaux (chapitre~15). Le travail à distance modifie le premier sans toucher au dernier~: un salarié qui travaille depuis un autre pays
peut créer des obligations fiscales et réglementaires pour l'employeur, ce qui est l'une des raisons pour lesquelles une entreprise dont le personnel traite des activités réglementées
peut exiger qu'il travaille depuis un bureau où elle est agréée.

\section{Tutoriel~: la même offre dans cinq villes}\label{tut:in:locations:tut}

\begin{tutorial}
\textbf{Objectif.} Calculer la \omterm{def:in:locations:pli}{rémunération disponible après logement} pour un package dans cinq villes, et les packages auxquels deux villes
sont équivalentes.
\textbf{État final~:} \cref{fig:in:locations:split,fig:in:locations:vslondon,fig:in:locations:pli}.
\begin{tutsteps}
\item \textbf{Villes.} \texttt{firm.locations.City} regroupe une localisation \texttt{firm.aftertax}, un loyer avec sa
mesure et sa période, un identifiant de registre~; \texttt{data/industry/rents.csv} contient les lignes de l'ONS et de HUD.
\item \textbf{Rémunération disponible.} \texttt{disposable(params, city, gross)} retranche le loyer annuel du salaire net
(\cref{lst:in:locations:disposable}).
\omcode{code/firm/locations/firm_locations.py}{35}{42}{Loyer annuel en dollars et rémunération disponible après logement.}\label{lst:in:locations:disposable}
\item \textbf{Comparaisons.} \texttt{equivalent} trouve le package d'une ville qui correspond à la rémunération disponible
d'une autre ville~; \texttt{crossing} calcule le package auquel deux villes sont équivalentes, par dichotomie.
\end{tutsteps}
\end{tutorial}

Pour 300\,000~\$ à Londres, le package new-yorkais équivalent est de 281\,047~\$, et pour 300\,000~\$ à New York, le
package londonien équivalent est de 318\,569~\$~; les deux villes se croisent vers 134\,100~\$ et 851\,900~\$.

\textbf{Pour aller plus loin.} Remplacer les statistiques de loyer par une mesure unique pour toutes les villes (un loyer médian de deux pièces
proposé à la location, issu d'un seul fournisseur, avec sa licence)~; ajouter le profil de report des primes (chapitre~13), qui modifie l'année d'imposition~;
ajouter la scolarité et les trajets domicile-travail pour une famille.

\section{Construction~: firm.locations}\label{bld:in:locations:locations}

\begin{build}
\bfield{Rôle} Comparer ce qu'un package laisse après impôt et logement dans différentes villes, et trouver les packages qui rendent
deux villes équivalentes.
\bfield{Interface} \texttt{firm.locations}~:
\texttt{City(name, location, rent\_month, currency, measure, period, source, visas)}~; \texttt{rent\_usd\_year}~; \texttt{disposable(params, city, gross)}~; \texttt{equivalent(params, a, gross\_a, b)}~;
\texttt{crossing(params, a, b, lo, hi)}~; sur \texttt{firm.aftertax}.
\bfield{Règles} Chaque loyer porte sa mesure, sa période et sa source~; les taux de change sont ceux des données fiscales~; un croisement
n'est signalé qu'à l'intérieur de la plage explorée, et l'appelant choisit la plage.
\bfield{Tests d'acceptation} \texttt{code/firm/locations/tests/}~: sans impôt, la rémunération disponible est le brut moins le loyer~;
\texttt{equivalent} inverse \texttt{disposable}~; deux villes qui ne se croisent jamais ne donnent aucun croisement.
\bfield{Pour aller plus loin} Ajustement au niveau de prix~; familles (deux revenus, enfants)~; une simulation de Monte Carlo sur la prime.
\end{build}

\begin{omsources}
\item ONS, Price Index of Private Rents~; HUD, Fair Market Rents FY2026~; City of Zurich, Mietpreiserhebung 2024~; Hong Kong
Rating and Valuation Department~; Eurostat, purchasing power parities.
\item USCIS, H-1B pages~; GOV.UK, Skilled Worker visa~; Singapore MOM, Employment Pass~; IND, required amounts 2026.
\item Données fiscales du chapitre~15 et \texttt{firm.aftertax}.
\end{omsources}

\section{Exercices}

\begin{exercise}[$\star$]\label{exo:in:locations:1}
Calculer le loyer annuel en dollars à Londres et à Hong Kong, d'après les statistiques du chapitre.
\end{exercise}

\begin{exercise}[$\star$]\label{exo:in:locations:2}
Calculer de combien la rémunération disponible à Hong Kong dépasse celle de Londres pour 300\,000~\$, et quelle part de l'écart représente l'impôt.
\end{exercise}

\begin{exercise}[$\star$]\label{exo:in:locations:3}
Calculer le salaire mensuel qu'un migrant hautement qualifié de 30 ans doit gagner aux Pays-Bas en 2026, et le montant annuel correspondant.
\end{exercise}

\begin{exercise}[$\star\star$]\label{exo:in:locations:4}
Expliquer pourquoi New York et Londres se croisent deux fois.
\end{exercise}

\begin{exercise}[$\star\star$]\label{exo:in:locations:5}
Le loyer de Zurich est un loyer médian net de charges pour 2024~; celui de Londres est une moyenne pour 2026. Expliquer dans
quel sens chaque écart biaise la comparaison.
\end{exercise}

\begin{exercise}[$\star\star$]\label{exo:in:locations:6}
Une candidate compare 300\,000~\$ à New York avec 280\,000~\$ à Chicago. Déterminer laquelle laisse le plus après impôt et
loyer, et préciser ce qu'elle devrait encore peser.
\end{exercise}

\begin{exercise}[$\star\star\star$]\label{exo:in:locations:7}
\emph{Programmation.} Doubler l'appartement de Hong Kong pour le porter à 110 mètres carrés. Que devient la rémunération disponible de Hong Kong à 300\,000~\$
et à son avance sur Zurich~?
\end{exercise}

\begin{exercise}[$\star\star\star$]\label{exo:in:locations:8}
\emph{Trouver la faille.} «~Zurich verse le plus après impôt et loyer parmi les villes européennes, c'est donc le meilleur
endroit où vivre avec le salaire d'un quant.~»
\end{exercise}

\section{Problème~: la même offre dans cinq villes}

\begin{problem}[{Problème du week-end --- la même offre dans cinq villes}]\label{pb:in:locations:1}
Un chercheur reçoit la même offre à 300\,000~\$ de la part des bureaux d'une même société à Londres, à New York, à Chicago, à Zurich et à Hong Kong.

\medskip\noindent\textbf{Partie I --- Les villes.}
\begin{enumerate}
\item Indiquer quelles villes de l'industrie le chapitre compare, et pourquoi celles-ci~?
\item Définir l'\omterm{def:in:locations:pli}{indice de niveau des prix} et la \omterm{def:in:locations:pli}{rémunération disponible après logement}.
\item Énoncer les cinq statistiques de loyer et expliquer en quoi elles diffèrent.
\item Expliquer pourquoi la colocation ne peut pas décider du lieu de travail.
\item Indiquer quelles mesures rendraient la comparaison plus cohérente.
\end{enumerate}

\medskip\noindent\textbf{Partie II --- Les chiffres.}
\begin{enumerate}[resume]
\item Énoncer la méthode.
\item Donner le salaire net, le loyer et la rémunération disponible dans chaque ville à 300\,000~\$.
\item Indiquer lequel, de l'impôt ou du loyer, sépare le plus les villes.
\item Donner les packages équivalents entre Londres et New York.
\item Déterminer à quels packages Londres et New York s'égalisent.
\end{enumerate}

\medskip\noindent\textbf{Partie III --- Visas et déménagement.}
\begin{enumerate}[resume]
\item Définir le \omterm{def:in:locations:visas}{visa pour profession spécialisée} et le \omterm{def:in:locations:visas}{visa de travailleur qualifié}.
\item Énoncer le plafond H-1B, la pondération à partir de l'exercice fiscal 2027 et le statut du versement de 100\,000~\$.
\item Énoncer les seuils britannique, singapourien et néerlandais.
\item Indiquer les trois coûts d'un déménagement au-delà de l'impôt et du loyer.
\item Expliquer pourquoi les sociétés de trading ont tendance à exiger une présence physique au bureau.
\end{enumerate}

\medskip\noindent\textbf{Partie IV --- Le verdict.}
\begin{enumerate}[resume]
\item Énoncer le \emph{résultat nommé}~: la rémunération disponible après impôt et après le loyer annuel d'un deux-chambres pour le package du chapitre dans
cinq villes, et le package brut qui égalise Londres et New York.
\item Indiquer comment les niveaux de prix modifieraient le \omterm{def:in:bank-markets-divisions:league}{classement}.
\item Indiquer ce qui changerait à 1~million de dollars.
\item Indiquer ce qu'ajoute la situation de visa pour un citoyen non américain qui choisit New York.
\item En deux phrases, conseiller le chercheur.
\end{enumerate}
\end{problem}

\section{Questions d'entretien}

\begin{interviewq}[$\star$ \iqroles{chercheur}]\label{iq:in:locations:1}
Comment compareriez-vous deux offres dans des pays différents~?
\end{interviewq}

\begin{interviewq}[$\star$ \iqroles{développeur}]\label{iq:in:locations:2}
Pourquoi une société pourrait-elle garder ses ingénieurs dans la ville de sa principale bourse plutôt que n'importe où~?
\end{interviewq}

\begin{interviewq}[$\star\star$ \iqroles{chercheur}]\label{iq:in:locations:3}
Les \omterm{def:in:pay-regulation-and-tax-by-location:rates}{taux moyens d'imposition} de deux villes sont égaux à 300\,000~\$. Leur \omterm{def:in:locations:pli}{rémunération disponible après logement} peut-elle
différer davantage que l'écart de loyer~?
\end{interviewq}

\begin{interviewq}[$\star\star$ \iqroles{chercheur, développeur}]\label{iq:in:locations:4}
Concevez une fonction qui renvoie le package d'équilibre entre deux villes, et dites ce qui peut mal tourner numériquement.
\end{interviewq}

\begin{interviewq}[$\star\star$ \iqroles{trader}]\label{iq:in:locations:5}
Pourquoi un teneur de marché pourrait-il garder des bureaux dans plusieurs fuseaux horaires plutôt qu'en un seul lieu~?
\end{interviewq}

\begin{interviewq}[$\star\star\star$ \iqroles{chercheur}]\label{iq:in:locations:6}
Le dollar recule de 10~\% face à la livre. Que devient le volet londonien de la comparaison du chapitre, et
que devient la rémunération d'un salarié londonien payé en livres~?
\end{interviewq}
